% Generated draft: requires mathematical and source review.
\documentclass[11pt]{article}
\usepackage[margin=25mm]{geometry}
\usepackage{fontspec}
\setmainfont{Latin Modern Roman}
\usepackage{amsmath,amssymb,mathtools,unicode-math}
\setmathfont{Latin Modern Math}
\usepackage{enumitem,needspace,xurl,hyperref}
\setlength{\parindent}{0pt}
\setlength{\parskip}{5pt}
\setlength{\emergencystretch}{3em}
\newcommand{\N}{\mathbb{N}}
\newcommand{\Z}{\mathbb{Z}}
\newcommand{\Q}{\mathbb{Q}}
\newcommand{\R}{\mathbb{R}}
\newcommand{\C}{\mathbb{C}}
\newcommand{\F}{\mathbb{F}}
\newcommand{\A}{\mathbb{A}}
\newcommand{\PP}{\mathbb{P}}
\newcommand{\E}{\mathbb{E}}
\newcommand{\eps}{\varepsilon}
\newcommand{\dd}{\,\mathrm{d}}
\begin{document}
\section*{UnsolvedMath ID: 700927}
\textbf{Source problem code:} catalogue entry
\subsection{Definitions and mathematical statement}
Let \(G\) be a group. It is \emph{periodic} if every element has finite order, and \emph{locally soluble} if every finitely generated subgroup is soluble. A group is soluble if its derived series reaches the trivial subgroup after finitely many steps. A group is nilpotent if it has a finite central series, and it is \emph{locally nilpotent} if every finitely generated subgroup is nilpotent.

A group \(K\) is \emph{Chernikov} if it has a normal subgroup of finite index isomorphic to
\[
\prod_{j=1}^{r} C_{p_j^\infty},
\]
where \(r\geq 0\) is finite, each \(p_j\) is a prime, and \(C_{p^\infty}\) is the quasicyclic \(p\)-group, isomorphic to the additive group \(\Z[1/p]/\Z\). The empty product is trivial, so finite groups are included. For \(H\leq G\), define
\[
C_G(H)=\{g\in G:gh=hg\text{ for every }h\in H\}.
\]

The question is whether, for every locally soluble periodic group \(G\), the existence of a finite nilpotent subgroup \(H\leq G\) such that \(C_G(H)\) is Chernikov implies the existence of a finite normal series
\[
\{1\}=G_0\leq G_1\leq\cdots\leq G_n=G
\]
whose factors \(G_{i+1}/G_i\) are locally nilpotent for every \(0\leq i<n\).

Here the standard invariant-series convention is used: every \(G_i\) is normal in \(G\). The supplied excerpt does not explicitly distinguish this convention from requiring only \(G_i\trianglelefteq G_{i+1}\); this terminology should be checked against the cited source. No uniform bound on \(n\) is requested.
\subsection{Short English statement}
Must every locally soluble periodic group containing a finite nilpotent subgroup with Chernikov centralizer have a finite normal series with locally nilpotent factors? The normal-series convention requires source review. The registry labels the problem open; that label is supplied metadata, not an independently verified assessment.
\subsection{Sources}
\begin{itemize}
\item Local UnsolvedMath registry, permanent ID \(700927\), source code \texttt{catalogue entry}, title \emph{Kourovka Notebook Problem 13.9}. The supplied record labels the problem open.
\item Supplied source attribution: \emph{The Kourovka Notebook (Unsolved Problems in Group Theory)}, 13th issue (1995), Problem 13.9; proposed by V. V. Belyaev and B. Hartley.
\item The supplied background also identifies the 21st issue edition (2026), \texttt{arXiv:1401.0300v48}. The collection description names E. I. Khukhro and V. D. Mazurov as editors. These bibliographic details are supplied attribution and have not been independently verified.
\end{itemize}
\end{document}
